\documentclass[11pt]{article}
\usepackage{fontspec}
\setmainfont{Times New Roman}
\setsansfont{Arial}
\setmonofont{Consolas}
\usepackage{amsmath,amssymb,mathtools}
\usepackage{geometry}
\usepackage{microtype}
\geometry{a4paper,margin=2.35cm}
\setlength{\parindent}{1.25em}
\setlength{\parskip}{0pt}
\makeatletter
\renewcommand\footnotesize{\@setfontsize\footnotesize{7.7}{8.7}\selectfont}
\makeatother
\emergencystretch=35em
\allowdisplaybreaks
\sloppy
\hbadness=10000
\vbadness=10000
\hfuzz=2pt
\vfuzz=10pt
\raggedbottom

\begin{document}

% Unit: Paper30 Section01 v001. Direct Interslavic Latin translation from the German final-audited source slice, source lines 137--389 / segments P30-S0014--P30-S0038. Footnote counter continues after Paper30 introduction.
\setcounter{footnote}{5}

\subsection*{\S~1. Teorija cělyh elementov}

Dano naj bude komutativno kolco\footnote{Kolco, kako jest znano, jest opreděljeno, kogda v množstvu dana jest relacija ravnosti, i takodže dvě operacije, složenje i množenje, ktore zadovoljajut obyčne zakony: složenje jest asociativno, komutativno i jednoznačno obratimo; množenje jest asociativno -- i komutativno, kogda jde o komutativno kolco -- i distributivno vzhodno k složenju. Ako osnovna definicija ravnosti ne jest množestvena identičnost, togda -- aby v vsakom podsistemu ostala ta sama definicija ravnosti -- takovy podsistem (podkolco, modul, ideal) mora vměstě s kakymkoli elementom vsebovati i vse jemu ravne elementy. Takodže pri vsakom razširjenju treba predpostaviti, že nova definicija ravnosti obhvača staru. Vsi pojmy kako «konečno mnogo elementov», «jednoznačno priredženje» i tako dalje razumějut se v smyslu definicije ravnosti; na primjer, «jest jeden i samo jeden element» znači, že jest samo jeden element do ravnosti. Vpročem, od ravnosti vsegda možno prějdti k množestvenoj identičnosti, sobravši ravne elementy v jeden klas i vvedši ty klasy kako nove elementy. Tut dane primětky o definiciji ravnosti pohodet od R. Hölzer.} $T$ bez nulovyh děliteljev, s jediničnym elementom $e$ množenja (aksiomy $III$ i $IV$); v $T$ vyznačeno jest fiksovano podkolco $R$, ktoro vsebuje jediničny element. Pojmy modul, porjadok i cělost vse vrěme odnoset se k tomu fiksovanomu kolcu $R$; pozdněje to radi kratkosti bude izpuščeno iz označenja.\footnote{Teorija cělyh elementov ostavaje validna, ako se dopustet nulove dělitelje i zato v $R$ predpostavi uslovje cěpov děliteljev; po §2 to takodže umožnja dokaz lemy. Poněže ta forma teorije cělyh elementov pozdněje ne bude upotrěbljena, na nju bude ukazano samo v primětkah. Vse definicije, o ktoryh jde rěč, ostavajut validne pri eksistenciji nulovyh děliteljev; večina iz njih, kako jest znano i lahko vidno, potrěbuje ještě slabše predpostavjenja. Ne ostavaje validna samo Dedekindova forma: «Čislo jest cělo, ako imaje oboločku». Bo pri pojavjenju nulovyh děliteljev kvocient dvoh konečnyh modulov ne mora ostati konečny; zato tut porjadok jest definovany neposrědnje, a ne kako modulny kvocient.} Elementy iz $R$ budut označane latinskymi bukvami, a elementy iz $T$ v obče grečskymi.

\textbf{1.} Sistem $M$ elementov iz $T$ nazyvaje se, kako obyčno, $R$-modulom, abo kratko modulom, ako vměstě s dvoma elementami $\alpha$ i $\beta$ vsebuje i jih razliku, a vměstě s $\alpha$ takodže $r\alpha$, gdje $r$ jest ljuby element iz $R$. Govori se, že $M$ jest dělimy črez $N$,
\[
        M\equiv 0 \pmod N,
\]
i že $M$ jest kratno, a $N$ dělitelj, ako vsaky element iz $M$ jest takodže vsebovany v $N$. $R$-moduly, ktoryh vsi elementy naležat k $R$, nazyvajut se idealy v $R$.

Očitno prěsěk -- najmenše obče kratno $[\ldots,M_i,\ldots]$ -- ljubogo množstva modulov iz $T$ jest znovo modul; takodže $T$ jest modul. Tako jednoznačno obstaja modul, vyvedeny iz ljubogo sistema $\Sigma$ elementov iz $T$, kako prěsěk vsih modulov, ktore vsebujejut $\Sigma$. Suma -- največji obči dělitelj $(\ldots,M_i,\ldots)$ -- ljubogo sistema modulov jest modul, vyvedeny iz jih sjedinenja. Produkt $MN$ dvoh modulov definovany jest kako modul, vyvedeny iz elementov $\mu\nu$, gdje $\mu$ proběga vse elementy iz $M$, a $\nu$ vse iz $N$. Produkt zato zadovolja asociativny i komutativny zakon, poněže ty zakony važe za elementy kolca. Iz distributivnogo zakona v $T$ dalje slěduje distributivny zakon za moduly:
\[
        (\ldots,M_i,\ldots)N=(\ldots,M_iN,\ldots),
\]
poněže po tom samom zakonu v $T$ prava i lěva strana sut modul, vyvedeny iz vsih elementov $\mu_i\nu$.

Poněže samo $R$ jest modul, uslovje, že $R$ jest oblast množiteljev, možno izraziti črez $RM\equiv0\pmod M$. A poněže po predpostavjenju $R$ vsebuje jediničny element, imajemo takodže $M\equiv0\pmod{RM}$, i zato $M=RM$.

Modul $M$ nazyvaje se konečny, to jest konečno porodženy, ako obstaja najmenej jeden sistem $\Sigma$ iz konečno mnogih elementov, iz ktorego $M$ jest vyvedeny; elementy $\mu_1,\ldots,\mu_n$ iz $\Sigma$ nazyvajut se modulna baza od
\[
        M=(\mu_1,\ldots,\mu_n).
\]
Suma i produkt dvoh -- i zato konečno mnogih -- konečnyh modulov sut znovo konečne, poněže sut vyvedene iz sjedinenja, odnosno iz produktov bazovyh elementov. Poslědnje slěduje iz predstavjenja $r_1\mu_1+\cdots+r_n\mu_n$ vsih elementov $\mu$ iz $M$ črez bazu i iz komutativnogo zakona množenja v $T$.

\textbf{2.} Razširjajuče kolco od $R$, ležeče v $T$ -- značit, kolco ktoro vsebuje vse elementy iz $R$ -- nazyvaje se $R$-porjadok, kratko porjadok. Prěsěk ljubogo množstva porjadkov iz $T$ jest znovo porjadok; takodže $T$ jest porjadok. Zato jednoznačno obstaja porjadok, vyvedeny iz ljubogo sistema $\Sigma$ elementov iz $T$. Sumy i produkty porjadkov definirajut se odgovarajuče kako za moduly; osoblivo $O^2=O$ za vsaky porjadok.

Vsaky porjadok jest Jednočasno $R$-modul; porjadok nazyvaje se konečny, ako jest konečny modul. Poněže sumy i produkty nastavajut tvorjenjem modulnoj sumy i produkta, po punktu 1 sumy i produkty konečnyh porjadkov sut znovo konečne porjadky. Dalej važi tranzitivny zakon: ako $A$ jest konečny $R$-porjadok, a $B$ konečny $A$-porjadok, togda $B$ jest takodže konečny $R$-porjadok. Bo $B$ Jednočasno stava $R$-porjadkom, znači $R$-modulom. Ako $\alpha_1,\ldots,\alpha_m$ tvoret modulnu bazu od $A$ vzhodno k $R$, a $\beta_1,\ldots,\beta_n$ taku bazu od $B$ vzhodno k $A$, togda produkty $\alpha_i\beta_j$ očitno tvoret modulnu bazu od $B$ vzhodno k $R$.

\textbf{3.} Element $\alpha$ iz $T$ nazyvaje se cěly vzhodno k $R$, kratko cěly, ako porjadok $R_\alpha$, vyvedeny iz $\alpha$, jest konečny; drugymi slovami, ako cěp děliteljev modulov
\[
        U_n=(\alpha^0,\alpha,\ldots,\alpha^{n-1})
\]
prěryvaje se po konečno mnogih krokah, $U_n=U_{n+1}=\cdots$.

Definicija dopušča takodže drugu formu, ktora bude upotrěbljena v punktu 4: element $\alpha$ iz $T$ nazyvaje se cěly, ako obstaja najmenej jeden glavny modul $C$, različny od nulovogo modula -- to jest $C$ jest vyvedeny iz elementa $\gamma\ne0$ -- takovy, že produkt porjadka $R_\alpha$ s $C$ jest konečny modul; drugymi slovami, že cěp modulov $CU_n$ prěryvaje se po konečno mnogih krokah.\footnote{Ako v $R$ dopustet se nulove dělitelje, treba zahtěvati eksistenciju regularnogo glavnogo modula; to jest $\gamma$ treba predpostaviti kako ne-nulovy dělitelj, regularny element, čto za kolca bez nulovyh děliteljev svodi se na $\gamma\ne0$.}

Nakonec definicija jest takodže identična s obyčnoj formoju: element $\alpha$ iz $T$ nazyvaje se cěly, ako zadovolja najmenej jedno jednačenje
\[
        \alpha^n+r_1\alpha^{n-1}+\cdots+r_n=0,
\]
gdje $r$ sut elementy iz $R$.

Različne definicije sut v samoj dějstvitelnosti identične. Poněže $R_\alpha$ sovpada s modulom, vyvedenym iz vsih stepenjev $\alpha$, za konečno $R_\alpha$ vsegda možno izbrati modulnu bazu iz tyh stepenjev. Ako taka baza jest dana napr. črez $e=\alpha^0,\alpha,\ldots,\alpha^{n-1}$, to znači, že cěp $U_n$ prěryvaje se pri $U_n$; obratno, iz $U_n=U_{n+1}=\cdots$ dobiva se ta baza. No prěryvanje cěpa $U_n$ jest takodže identično s eksistencijo jednačenja
\[
        \alpha^n+r_1\alpha^{n-1}+\cdots+r_n=0.
\]
S $R_\alpha$ jest takodže $CR_\alpha$ konečny, zato cěp $CU_n$ se prěryvaje; obratno, iz konečnosti $CR_\alpha$, to jest iz prěryvanja $CU_n$, slěduje takodže prěryvanje $U_n$ i s tym konečnost $R_\alpha$. Bo poněže $C$ jest predpostavjen kako glavny modul različny od nulovogo modula, $CU_n=CU_{n+1}$ dava takodže $U_n=U_{n+1}$.

Kolcevo svojstvo cělyh elementov i tranzitivny zakon cěloj zavisimosti opirajut se na slědujuču lemu.

\medskip
\noindent\textbf{Lema.} Ako $O$ jest konečny porjadok iz $T$, a $\alpha$ ljuby element iz $O$, togda porjadok $R_\alpha$, vyvedeny iz $\alpha$, jest konečny; drugymi slovami, vsi elementy konečnogo porjadka sut cěle.

Dokaz dobiva se po obyčnyh zaključkah. Ako $\omega_1,\ldots,\omega_n$ jest modulna baza od $O$, togda zaradi kolcevogo svojstva takodže $\alpha\omega_i$ naleži k $O$. Zato nastava
\[
        \alpha\omega_i=r_{i1}\omega_1+\cdots+r_{in}\omega_n,
\]
i iz togo
\[
        \det(r_{ij}-e_{ij}\alpha)=0,
\]
s $e_{ij}=0$ za $i\ne j$, $e_{ii}=e$, čim $\alpha$ jest dokazany kako cěly. Za nulovost determinanta upotrěblja se, že $T$, i zato $O$, jest predpostavjeno bez nulovyh děliteljev.\footnote{Ako v $R$ jest predpostavjeno uslovje cěpov děliteljev, ale dopustet se nulove dělitelje, togda lema jest neposrědnji poslědok modulnogo teorema v §2, po ktorom cěpovo uslovje prěnosi se na moduly iz $O$. I toj dokaz za čislovo polje pochodi od Dedekind (4. izd., §173, II) i jest prva upotrěba cěpovyh uslovij v literaturi. V poslědkah, ktore budut dane pod punktom 4, Dedekind vměsto togo upotrěblja ještě specialnějši fakt, že čislo klasov ostatkov modulo ideal v algebraičnom čislovom polju jest konečno.}

Sistem $G$ vsih cělyh elementov iz $T$ tvori porjadok. $G$ vsebuje $R$, bo elementy iz $R$ sut cěle, poněže $R$ tvori konečny $R$-modul s jediničnym elementom kako bazoju. Ale $G$ imaje takodže kolcevo svojstvo: porjadok $R_{\alpha,\beta}$, vyvedeny iz dvoh ljubych cělyh elementov $\alpha,\beta$, stava ravny produktu $R_\alpha R_\beta$, i zato jest konečny. Po lemi, $\alpha+\beta$ i $\alpha\beta$ sut zato cěle kako elementy konečnogo porjadka.

\medskip
\noindent\textbf{Tranzitivny zakon cěloj zavisimosti.} Ako $O$ jest porjadok sostojajuči iz cělyh elementov, togda vsaky element iz $T$, cěly vzhodno k $O$, jest takodže cěly vzhodno k $R$. Po definiciji $\alpha$ zadovolja jednačenje
\[
        \alpha^m+\omega_1\alpha^{m-1}+\cdots+\omega_m=0,
\]
i zato jest takodže cěly vzhodno k porjadku $O'=R_{\omega_1,\ldots,\omega_m}$, vyvedenomu iz $\omega_1,\ldots,\omega_m$, ktory jest konečny kako produkt $R_{\omega_1}\cdots R_{\omega_m}$. Po tranzitivnom zakonu danom v punktu 2, konečny $O'$-porjadok $O'_\alpha$, vyvedeny iz $\alpha$, stava takodže konečny $R$-porjadok, i $\alpha$ jest cěly po lemi kako element konečnogo porjadka.

\textbf{4.} Kolco $R$ nazyvaje se integralno zamknjeno v $T$, ako vsaky element iz $T$, cěly vzhodno k $R$, naleži k $R$. Po drugoj formě definicije cělyh elementov v punktu 3, element $\alpha$ iz $T$ naleži k $R$ togda i samo togda, kogda obstaja najmenej jeden glavny modul $C$, različny od nulovogo modula, takovy že $CR_\alpha$ tvori konečny modul.\footnote{Ako v $R$ dopustet se nulove dělitelje, togda $C$ znovo treba predpostaviti kako regularny glavny modul; takodže element $c$ v poslědku $I$ treba predpostaviti kako regularny.}

Iz pojma integralno zamknjenogo kolca Dedekind (3. izd., §172) vyvel dva važne poslědky, ktore dozvaljajut upotrěbiti toj pojem za teoriju idealov:

\medskip
\noindent\textbf{Dedekindov poslědok $I$.} $R$ naj bude integralno zamknjeno v $T$, i takodže naj v $R$ kako daljnja predpostavka važi uslovje cěpov děliteljev za idealy (aksioma $I$). Ako $\beta$ jest necěly element iz $T$, a $c\ne0$ element iz $R$, togda jest eksponent $\rho>0$ takovy, že v redu elementov
\[
        c,\;c\beta,\ldots,c\beta^\nu,\ldots
\]
elementy $c,\ldots,c\beta^\rho$ sut cěle, a vsi ostale necěle.

Ako by vsi elementy $c\beta^\nu$ byli cěle, togda by oni zaradi integralnoj zamknjenosti $R$ naležali k $R$. Togda by cěp modulov $C,CB_1,\ldots,CB_\nu,\ldots$ -- gdje $C$ označa modul vyvedeny iz $c$, a $B_\nu$ modul vyvedeny iz $1,\beta,\ldots,\beta^\nu$ -- prěšla v cěp idealov $C_\nu=(c,c\beta,\ldots,c\beta^\nu)$, ktora po predpostavjenju prěryvaje se po konečno mnogih krokah. No togda by, protivno predpostavjenju, $R_\beta$ bylo konečno i $\beta$ cěly; zato medžu elementami $c\beta^\nu$ sut necěle. Dalej, vměstě s kakymkoli necělym elementom vse slědujuče sut necěle: bo ako $c\beta^r$ jest cěly, znači v $R$, togda $c\beta^\rho$ za $\rho<r$ zadovolja jednačenje
\[
        (c\beta^\rho)^r-c^{\,r-\rho}(c\beta^r)^\rho=0
\]
i zato jest takodže cěly. Ako zato $c\beta^{\rho+1}$ jest prvi necěly element, togda -- poněže $c$ jest cěly -- $\rho>0$ i jest iskany eksponent.

Ako specializiramo razširjajuče kolco $T$ do polja kvocientov od $R$ -- to jest do polja, ktoro nastava adjunkcijo vsih parov elementov\footnote{V znanoj formulaciji Steinitz, J. f. M. 137. Ako v $R$ dopustet se nulove dělitelje, vměsto polja kvocientov mora nastupiti kolco kvocientov, črez adjunkciju vsih kvocientnyh parov, v ktoryh imenitelj jest regularny element iz $R$. Tut Dedekindov poslědok $II$ uže ne jest ispolnjen, bo $n$ v obče ne bude regularny element.} -- togda iz togo slěduje:

\medskip
\noindent\textbf{Dedekindov poslědok $II$.} $R$ naj bude integralno zamknjeno v svojem polju kvocientov $K$, i v $R$ naj važi uslovje cěpov děliteljev za idealy. Ako $\beta$ jest necěly element iz $K$, togda $\beta$ dopusča predstavjenje kako kvocient elementov iz $R$,
\[
        \beta=m/n,
\]
takovo, že takodže $m^2/n$ jest necěly.

Postavimo $\beta=b/c$ s $c\ne0$. Togda v redu $c,c\beta=b,c\beta^2,\ldots$ prva dva elementa sigurno sut cěle, zato $\rho>1$, gdje $\rho$ označa eksponent iz poslědka $I$. Ako postavimo
\[
        m=c\beta^\rho,\qquad n=c\beta^{\rho-1},
\]
togda $m/n=\beta$ i $m^2/n=c\beta^{\rho+1}$ jest necěly.

Tranzitivny zakon integralnoj zamknjenosti imaje slědujuču formu: ako $R$ jest ljubo kolco iz $T$, i $G$ označa sistem vsih elementov iz $T$, cělyh vzhodno k $R$, togda $G$ sostojit takodže iz vsih elementov iz $T$, cělyh vzhodno k $G$. Bo vsaky element, cěly vzhodno k $G$, jest takodže cěly vzhodno k $R$ po tranzitivnom zakonu v punktu 3, i zato naleži k $G$.

\end{document}
